Blockchainové technológie predstavujú významnú oblasť súčasnej
informatiky a finančných technológií. Blockchain možno charakterizovať
ako distribuovaný systém, v ktorom sú údaje o transakciách ukladané
v chronologicky usporiadaných blokoch. Verejný charakter blockchainu
umožňuje používateľom overovať transakcie bez centrálneho správcu,
pričom zároveň vytvára rozsiahly zdroj údajov vhodný na analytické
spracovanie.

Analýza blockchainových transakcií je dôležitá najmä z hľadiska
identifikácie podozrivých aktivít, prania špinavých peňazí a ďalších
nelegálnych činností. Na tento účel sa využívajú metódy strojového
učenia, grafové modely a algoritmy na detekciu anomálií. Transakcie
možno reprezentovať ako graf, v ktorom vrcholy predstavujú transakcie
a hrany vyjadrujú tok prostriedkov medzi nimi.

Na tomto momente vsetko dole nie je finalna cast lebo zatial neviem co presne budem robit

Predmetom tejto práce je verejne dostupný anonymizovaný Elliptic
Bitcoin Dataset. Dataset obsahuje údaje o bitcoinových transakciách,
ich vzájomných vzťahoch, časových krokoch a ďalších vlastnostiach.
Keďže identifikátory transakcií v datasete sú anonymizované, nie je
možné automaticky predpokladať, že zodpovedajú priamo identifikátorom
transakcií vo verejnom blockchaine.

Cieľom bakalárskej práce je preskúmať možnosti spätného dohľadania
transakcií z anonymizovaného Elliptic datasetu vo verejnom Bitcoin
blockchaine. Práca sa zameria na analýzu dostupných identifikátorov
a vlastností transakcií, návrh metódy vyhľadávania možných zhôd
a overenie správnosti nájdených výsledkov.

Na dosiahnutie hlavného cieľa budú riešené tieto úlohy:

\begin{itemize}
    \item analyzovať štruktúru a obsah Elliptic Bitcoin Datasetu;
    \item preskúmať možnosti identifikácie transakcií vo verejnom blockchaine;
    \item navrhnúť metódu párovania anonymizovaných údajov s blockchainovými údajmi;
    \item implementovať navrhnutý postup;
    \item experimentálne overiť úspešnosť a obmedzenia navrhnutej metódy.
\end{itemize}

Výsledkom práce bude vyhodnotenie, do akej miery je možné spojiť
transakcie z anonymizovaného datasetu s reálnymi transakciami vo
verejnom blockchaine. Súčasťou práce bude aj diskusia o presnosti
výsledkov, použitých vlastnostiach transakcií a obmedzeniach
deanonymizácie blockchainových údajov.

Práca je členená do niekoľkých častí. Prvá časť opisuje blockchain,
Bitcoinové transakcie a vlastnosti použitého datasetu. Ďalšia časť
sa venuje návrhu metódy vyhľadávania a párovania transakcií.
Nasledujúca časť opisuje implementáciu a experimentálne vyhodnotenie.
V závere sú zhrnuté dosiahnuté výsledky a uvedené možnosti ďalšieho
výskumu.